\documentclass[11pt]{article}
\usepackage[margin=1in]{geometry}
\usepackage[T1]{fontenc}
\usepackage{lmodern}
\usepackage{amsmath,amssymb,amsthm,mathtools}
\usepackage[hidelinks]{hyperref}

\newtheorem{theorem}{Theorem}
\newtheorem{proposition}{Proposition}
\newtheorem{remark}{Remark}
\DeclareMathOperator{\esssup}{ess\,sup}
\newcommand{\R}{\mathbb R}
\newcommand{\F}{\mathcal F}
\newcommand{\ip}[2]{\langle #1,#2\rangle}
\newcommand{\norm}[1]{\lVert #1\rVert}

\title{E58: Sharp Bound for the Band-Selective Green Multiplier}
\author{Richard J. Reyes\\Independent Researcher}
\date{October 26, 2026}

\begin{document}
\maketitle

\begin{abstract}
For the fourth-order band-selective operator
$\mathcal L=r+a(\Delta+k_\star^2)^2$ on $\R^d$, with $r>0$, $a\geq0$,
and $k_\star\geq0$, we prove the pointwise Fourier-multiplier bound
$0<G(k)\leq 1/r$. For $a>0$, the bound is sharp exactly on the selected
shell $|k|=k_\star$. We further prove the exact $L^2$ resolvent norm,
identify the discrete-spectrum qualification, and distinguish positivity
of the Fourier multiplier from positivity of the physical-space Green
function.
\end{abstract}

\section{Definitions and principal result}

Fix an integer $d\geq1$. On the Schwartz class $\mathcal S(\R^d)$,
consider
\begin{equation}\label{eq:L}
  \mathcal L := r+a(\Delta+k_\star^2)^2,
  \qquad r>0,\quad a\geq0,\quad k_\star\geq0.
\end{equation}
We use the Fourier-transform convention
\begin{equation}\label{eq:fourier}
  \widehat u(k)=\int_{\R^d}e^{-ik\cdot x}u(x)\,dx,
  \qquad
  u(x)=\frac{1}{(2\pi)^d}\int_{\R^d}e^{ik\cdot x}\widehat u(k)\,dk.
\end{equation}
Write
\begin{equation}\label{eq:symbol}
  \ell(k):=r+a\bigl(|k|^2-k_\star^2\bigr)^2,
  \qquad
  G(k):=\frac{1}{\ell(k)}.
\end{equation}

\begin{theorem}[E58: positive, sharply bounded Green multiplier]
\label{thm:E58}
For every $k\in\R^d$,
\begin{equation}\label{eq:mainbound}
  \boxed{\quad 0<G(k)\leq\frac1r.\quad}
\end{equation}
If $a>0$, equality holds if and only if $|k|=k_\star$; in particular,
\begin{equation}\label{eq:sup}
  \sup_{k\in\R^d}G(k)=\frac1r.
\end{equation}
If $a=0$, $G(k)=1/r$ for all $k$.
\end{theorem}

\begin{proof}
Since $\widehat{\Delta u}(k)=-|k|^2\widehat u(k)$,
\begin{align*}
  \widehat{(\Delta+k_\star^2)u}(k)
    &=(k_\star^2-|k|^2)\widehat u(k),\\
  \widehat{(\Delta+k_\star^2)^2u}(k)
    &=(|k|^2-k_\star^2)^2\widehat u(k).
\end{align*}
Therefore $\widehat{\mathcal Lu}(k)=\ell(k)\widehat u(k)$ with
$\ell$ as in \eqref{eq:symbol}. For every $k$,
\begin{equation}\label{eq:denom}
  \bigl(|k|^2-k_\star^2\bigr)^2\geq0,
  \qquad
  a\bigl(|k|^2-k_\star^2\bigr)^2\geq0,
\end{equation}
and hence
\begin{equation}
  0<r\leq r+a\bigl(|k|^2-k_\star^2\bigr)^2=\ell(k).
\end{equation}
Taking reciprocals of positive numbers reverses the inequality:
\begin{equation}
  0<\frac{1}{\ell(k)}\leq\frac1r,
\end{equation}
which proves \eqref{eq:mainbound}.

For $a>0$, equality holds precisely when
$a(|k|^2-k_\star^2)^2=0$, or equivalently when
$|k|^2=k_\star^2$. Since $k_\star\geq0$, this is exactly
$|k|=k_\star$. That set is nonempty in $\R^d$, even if
$k_\star=0$, proving \eqref{eq:sup}. When $a=0$, the symbol is
identically $r$, so $G\equiv1/r$.
\end{proof}

\section{Inverse operator and exact $L^2$ estimate}

\begin{proposition}[Sharp continuum resolvent bound]\label{prop:L2}
Let $\mathcal L$ be as in \eqref{eq:L}. For each $f\in L^2(\R^d)$,
there exists a unique $u\in L^2(\R^d)$ satisfying
$\mathcal Lu=f$ in the sense of tempered distributions. It is defined
by
\begin{equation}\label{eq:solution}
  \widehat u(k)=G(k)\widehat f(k).
\end{equation}
Moreover,
\begin{equation}\label{eq:operatornorm}
  \boxed{\quad
  \norm{\mathcal L^{-1}}_{L^2\to L^2}=\frac1r.
  \quad}
\end{equation}
If $a>0$, then $u\in H^4(\R^d)$.
\end{proposition}

\begin{proof}
By Theorem~\ref{thm:E58}, $G\in L^\infty(\R^d)$ and
$\norm{G}_{L^\infty}\leq1/r$. For any $f\in L^2$, the product
$G\widehat f\in L^2$, so Plancherel's theorem defines the unique
$u\in L^2$ with \eqref{eq:solution} and gives
\begin{equation}\label{eq:upper}
  \norm{u}_{2}^2
  =\frac{1}{(2\pi)^d}\int_{\R^d}|G(k)|^2
   |\widehat f(k)|^2\,dk
  \leq\frac{1}{r^2}\norm{f}_2^2.
\end{equation}
Since $\ell(k)G(k)=1$, equation \eqref{eq:solution} solves
$\mathcal Lu=f$ distributionally. Conversely, if
$\mathcal Lu=0$ and $u\in L^2$, then
$\ell\widehat u=0$ and $\ell>0$, hence $u=0$.

To show that \eqref{eq:upper} is sharp, fix $\varepsilon>0$ and let
\begin{equation}
  E_\varepsilon
  :=\left\{k\in\R^d:
      \bigl||k|^2-k_\star^2\bigr|<\varepsilon\right\}.
\end{equation}
This measurable set has positive, finite measure for every
$\varepsilon>0$ (including when $k_\star=0$). Choose a nonzero
$f_\varepsilon\in L^2$ whose Fourier transform is supported in
$E_\varepsilon$. On this set,
\begin{equation}
  G(k)\geq\frac{1}{r+a\varepsilon^2},
\end{equation}
so Plancherel implies
\begin{equation}
  \frac{\norm{\mathcal L^{-1}f_\varepsilon}_2}
       {\norm{f_\varepsilon}_2}
  \geq\frac{1}{r+a\varepsilon^2}.
\end{equation}
Letting $\varepsilon\downarrow0$ gives
$\norm{\mathcal L^{-1}}_{L^2\to L^2}\geq1/r$,
which, together with \eqref{eq:upper}, proves
\eqref{eq:operatornorm}.

Finally, if $a>0$, the function
$(1+|k|^2)^2G(k)$ is bounded: its denominator grows like
$a|k|^4$ at high frequency and stays strictly positive at finite
frequency. Hence
\begin{equation}
  \int_{\R^d}(1+|k|^2)^4|\widehat u(k)|^2\,dk
  \leq C^2\int_{\R^d}|\widehat f(k)|^2\,dk<\infty,
\end{equation}
so $u\in H^4(\R^d)$.
\end{proof}

\begin{remark}[Norm attainment]
For $a>0$, the equality set $\{k:|k|=k_\star\}$ has
Lebesgue measure zero. Thus the $L^2\to L^2$ operator norm in
\eqref{eq:operatornorm} is approached by concentrating Fourier mass
near the shell, but is not attained by any nonzero $L^2$ input.
For $a=0$, every nonzero input attains it.
\end{remark}

\section{Energy (coercivity) proof of the upper estimate}

When $a>0$ and $u\in H^4(\R^d)$, integration by parts, or
Plancherel's theorem, yields
\begin{equation}\label{eq:coercivity}
\begin{aligned}
  \ip{u}{\mathcal Lu}_{L^2}
  &=r\norm{u}_2^2
    +a\norm{(\Delta+k_\star^2)u}_2^2\\
  &\geq r\norm{u}_2^2.
\end{aligned}
\end{equation}
If $\mathcal Lu=f$, then Cauchy--Schwarz gives
\begin{equation}
  r\norm{u}_2^2
  \leq\operatorname{Re}\ip{u}{f}_{L^2}
  \leq\norm{u}_2\norm{f}_2.
\end{equation}
Dividing by $\norm{u}_2$ when it is nonzero proves
$\norm{u}_2\leq r^{-1}\norm{f}_2$; when $u=0$, the estimate is
immediate. For $a=0$, $\mathcal L=rI$ and the estimate is an
identity.

\section{Band selectivity and the discrete-spectrum caveat}

Assume $a>0$ and set
$\delta:=\bigl||k|^2-k_\star^2\bigr|\geq0$. Then
\begin{equation}
  G(\delta)=\frac{1}{r+a\delta^2},
  \qquad
  \frac{dG}{d\delta}
  =-\frac{2a\delta}{(r+a\delta^2)^2}<0
  \quad(\delta>0).
\end{equation}
Thus the response is strictly largest on the shell and strictly
decreases with detuning. At high detuning,
\begin{equation}
  G(\delta)\sim\frac{1}{a\delta^2}
  \qquad(\delta\to\infty).
\end{equation}

For a periodic box, the same Fourier-multiplier inequality holds
at every admissible discrete wavevector $k\in\Lambda$. However,
\begin{equation}
  \norm{\mathcal L^{-1}}_{L^2\to L^2}
  =\sup_{k\in\Lambda}G(k)
  \leq\frac1r,
\end{equation}
with equality for $a>0$ if and only if the discrete spectrum
contains a wavevector with $|k|=k_\star$. The maximum is smaller
than $1/r$ otherwise.

\section{Fourier multiplier versus physical-space Green function}

The inequality $0<G(k)\leq1/r$ is a \emph{Fourier-space}
statement. It does not imply that the inverse Fourier transform
$g=\F^{-1}G$ is pointwise nonnegative, nor does it provide a
pointwise bound $g(x)\leq1/r$.

To see why nonnegativity fails, consider $d=1$, $k_\star=0$,
$r>0$, $a>0$. Write
\begin{equation}
  b:=\left(\frac ra\right)^{1/4},
  \qquad t:=\frac{b|x|}{\sqrt2}.
\end{equation}
The explicit spatial Green function for
$\mathcal L=r+a\,d^4/dx^4$ is
\begin{equation}\label{eq:physicalgreen}
  g(x)=\frac{1}{2\sqrt2\,a^{1/4}r^{3/4}}
       e^{-t}\bigl(\cos t+\sin t\bigr).
\end{equation}
For completeness, the formula follows from the even, decaying
solution to $ag''''+rg=0$ away from $x=0$. Indeed, with
$c=b/\sqrt2$, the function
$h(x)=e^{-c|x|}(\cos(c|x|)+\sin(c|x|))$
satisfies $ah''''+rh=0$ for $x\ne0$, is continuous with
$h'(0^+)=h'(0^-)=0$, and has
$h'''(0^+)-h'''(0^-)=8c^3$. Multiplying by
$C=(8ac^3)^{-1}=(2\sqrt2\,a^{1/4}r^{3/4})^{-1}$ gives
$ag''''+rg=\delta_0$ in the sense of distributions.
Since the Fourier transform of this identity yields
$\widehat g(k)=1/(r+ak^4)$, equation
\eqref{eq:physicalgreen} is exactly the inverse Fourier transform
of $G$ in this example.

At $t=\pi$, one has
$\cos t+\sin t=-1$, so $g(x)<0$, even though
$G(k)>0$ for every $k$. Consequently, positivity of the
Fourier multiplier must not be confused with pointwise positivity
of the spatial Green function.

\section{Scope and conclusion}

The assumptions $r>0$ and $a\geq0$ are essential to the
unconditional bound. If $r=0$ and $a>0$, the Fourier symbol
vanishes on the selected shell. If $a<0$, it is not globally
bounded below by $r$ and can vanish at real frequencies. The
result establishes a bounded \emph{linear} spectral response;
it does not, by itself, establish nonlinear pattern stability,
spatial localization, or physical wavelength selection.

\begin{equation*}
  \boxed{
    \forall k\in\R^d:\quad
    0<\frac{1}{r+a(|k|^2-k_\star^2)^2}\leq\frac1r,
    \qquad
    \norm{\mathcal L^{-1}}_{L^2\to L^2}=\frac1r.
  }
\end{equation*}

\end{document}
